\section{Five and six rows}\label{sec:five}

\subsection{Boundary padding}

For a blackout $S$ of the $h\times m$ grid and $L,R\ge0$, form $S_{L,R}$ by
adding columns $1-L,\dots,0$ and $m+1,\dots,m+R$, black only in a chosen
\emph{padding row}, row $1$ or row $h$.

\begin{theorem}[finite collar]\label{thm:collar}
Let $h\ge3$, fix the padding row and the padded sides, and let $L^*$ and $R^*$ be $W_h$ of
Lemma~\ref{lem:span} on padded sides and $0$ otherwise. Then $S_{L,R}$ is valid for all
$L,R\ge0$ vanishing on unpadded sides if and only if $S_{L^*,R^*}$ is valid.
\end{theorem}

\begin{proof}
Set $S^*=S_{L^*,R^*}$. Necessity follows by this choice. Suppose $S^*$ is valid
and put $U=S_{L,R}$.

A tilted rectangle has at most one corner in row $1$ and at most one
in row $h$ (Lemma~\ref{lem:leftmost} and~\eqref{eq:tilt-param}).
A weak tilted rectangle of $U$ therefore has a black corner in an
original column $e$. Its span is bounded by Lemma~\ref{lem:span}, so
its columns lie in $[e-W_h,e+W_h]\cap[1-L,m+R]\subseteq[1-L^*,m+R^*]$.
It lies in $S^*$ and has the same presentation there.

Suppose distinct $R_1,R_2$ in $U$ show the same $P$.

\emph{Both axis-aligned.} Every added column has degree one in $G_U$,
so any cycle already lies in $G_{S^*}$. Since $S^*$ is valid and has
at least three columns, Lemma~\ref{lem:girth} makes $U$ axis-valid.

\emph{Both tilted.} Lemma~\ref{lem:threecorners} gives $|P|\le2$.
Both rectangles are weak, so both lie in $S^*$ with presentation $P$,
contradicting its validity.

\emph{One of each.} Take $R_1$ tilted, hence already in $S^*$.
A column of $R_2$ meeting $P$ also belongs to $R_1$. A column not
meeting $P$ has two black corners, so it must be original.
Thus $R_2$ also lies in $S^*$ with the same presentation, again a
contradiction.
\end{proof}

\subsection{Five rows}

For the constructions, list the black rows in each column from left
to right. Start with the blackout $S_8$ given by $14,\allowbreak 45,\allowbreak 25,\allowbreak 5,\allowbreak 5,\allowbreak 34,\allowbreak 23,\allowbreak 12$ of the $5\times8$ grid, so that $B_1=\{1,4\}$ and $B_4=\{5\}$.

\begin{figure}[!t]
\centering
\begin{tikzpicture}[x=0.8cm,y=0.8cm,baseline=(current bounding box.north)]
  \begin{scope}[xshift=0.000cm]
    \draw[black,line width=0.32pt,dash pattern=on 1.20pt off 1.20pt] (1,-1) -- (13,-1) (1,-2) -- (13,-2) (1,-3) -- (13,-3) (1,-4) -- (13,-4) (1,-5) -- (13,-5) (1,-1) -- (1,-5) (2,-1) -- (2,-5) (3,-1) -- (3,-5) (4,-1) -- (4,-5) (5,-1) -- (5,-5) (6,-1) -- (6,-5) (7,-1) -- (7,-5) (8,-1) -- (8,-5) (9,-1) -- (9,-5) (10,-1) -- (10,-5) (11,-1) -- (11,-5) (12,-1) -- (12,-5) (13,-1) -- (13,-5);
    \filldraw[fill=kgray,draw=black,line width=0.48pt] (1,-1) circle (0.3);
    \draw[black,line width=0.48pt] (0.788,-0.788) -- (1.212,-1.212) (0.788,-1.212) -- (1.212,-0.788);
    \filldraw[fill=white,draw=black,line width=0.48pt] (2,-1) circle (0.3);
    \filldraw[fill=white,draw=black,line width=0.48pt] (3,-1) circle (0.3);
    \filldraw[fill=white,draw=black,line width=0.48pt] (4,-1) circle (0.3);
    \filldraw[fill=white,draw=black,line width=0.48pt] (5,-1) circle (0.3);
    \filldraw[fill=white,draw=black,line width=0.48pt] (6,-1) circle (0.3);
    \filldraw[fill=white,draw=black,line width=0.48pt] (7,-1) circle (0.3);
    \filldraw[fill=kgray,draw=black,line width=0.48pt] (8,-1) circle (0.3);
    \draw[black,line width=0.48pt] (7.788,-0.788) -- (8.212,-1.212) (7.788,-1.212) -- (8.212,-0.788);
    \filldraw[fill=kgray,draw=black,line width=0.48pt] (9,-1) circle (0.3);
    \draw[black,line width=0.48pt] (8.788,-0.788) -- (9.212,-1.212) (8.788,-1.212) -- (9.212,-0.788);
    \filldraw[fill=kgray,draw=black,line width=0.48pt] (10,-1) circle (0.3);
    \draw[black,line width=0.48pt] (9.788,-0.788) -- (10.212,-1.212) (9.788,-1.212) -- (10.212,-0.788);
    \filldraw[fill=kgray,draw=black,line width=0.48pt] (11,-1) circle (0.3);
    \draw[black,line width=0.48pt] (10.788,-0.788) -- (11.212,-1.212) (10.788,-1.212) -- (11.212,-0.788);
    \filldraw[fill=kgray,draw=black,line width=0.48pt] (12,-1) circle (0.3);
    \draw[black,line width=0.48pt] (11.788,-0.788) -- (12.212,-1.212) (11.788,-1.212) -- (12.212,-0.788);
    \filldraw[fill=kgray,draw=black,line width=0.48pt] (13,-1) circle (0.3);
    \draw[black,line width=0.48pt] (12.788,-0.788) -- (13.212,-1.212) (12.788,-1.212) -- (13.212,-0.788);
    \filldraw[fill=white,draw=black,line width=0.48pt] (1,-2) circle (0.3);
    \filldraw[fill=white,draw=black,line width=0.48pt] (2,-2) circle (0.3);
    \filldraw[fill=kgray,draw=black,line width=0.48pt] (3,-2) circle (0.3);
    \draw[black,line width=0.48pt] (2.788,-1.788) -- (3.212,-2.212) (2.788,-2.212) -- (3.212,-1.788);
    \filldraw[fill=white,draw=black,line width=0.48pt] (4,-2) circle (0.3);
    \filldraw[fill=white,draw=black,line width=0.48pt] (5,-2) circle (0.3);
    \filldraw[fill=white,draw=black,line width=0.48pt] (6,-2) circle (0.3);
    \filldraw[fill=kgray,draw=black,line width=0.48pt] (7,-2) circle (0.3);
    \draw[black,line width=0.48pt] (6.788,-1.788) -- (7.212,-2.212) (6.788,-2.212) -- (7.212,-1.788);
    \filldraw[fill=kgray,draw=black,line width=0.48pt] (8,-2) circle (0.3);
    \draw[black,line width=0.48pt] (7.788,-1.788) -- (8.212,-2.212) (7.788,-2.212) -- (8.212,-1.788);
    \filldraw[fill=white,draw=black,line width=0.48pt] (9,-2) circle (0.3);
    \filldraw[fill=white,draw=black,line width=0.48pt] (10,-2) circle (0.3);
    \filldraw[fill=white,draw=black,line width=0.48pt] (11,-2) circle (0.3);
    \filldraw[fill=white,draw=black,line width=0.48pt] (12,-2) circle (0.3);
    \filldraw[fill=white,draw=black,line width=0.48pt] (13,-2) circle (0.3);
    \filldraw[fill=white,draw=black,line width=0.48pt] (1,-3) circle (0.3);
    \filldraw[fill=white,draw=black,line width=0.48pt] (2,-3) circle (0.3);
    \filldraw[fill=white,draw=black,line width=0.48pt] (3,-3) circle (0.3);
    \filldraw[fill=white,draw=black,line width=0.48pt] (4,-3) circle (0.3);
    \filldraw[fill=white,draw=black,line width=0.48pt] (5,-3) circle (0.3);
    \filldraw[fill=kgray,draw=black,line width=0.48pt] (6,-3) circle (0.3);
    \draw[black,line width=0.48pt] (5.788,-2.788) -- (6.212,-3.212) (5.788,-3.212) -- (6.212,-2.788);
    \filldraw[fill=kgray,draw=black,line width=0.48pt] (7,-3) circle (0.3);
    \draw[black,line width=0.48pt] (6.788,-2.788) -- (7.212,-3.212) (6.788,-3.212) -- (7.212,-2.788);
    \filldraw[fill=white,draw=black,line width=0.48pt] (8,-3) circle (0.3);
    \filldraw[fill=white,draw=black,line width=0.48pt] (9,-3) circle (0.3);
    \filldraw[fill=white,draw=black,line width=0.48pt] (10,-3) circle (0.3);
    \filldraw[fill=white,draw=black,line width=0.48pt] (11,-3) circle (0.3);
    \filldraw[fill=white,draw=black,line width=0.48pt] (12,-3) circle (0.3);
    \filldraw[fill=white,draw=black,line width=0.48pt] (13,-3) circle (0.3);
    \filldraw[fill=kgray,draw=black,line width=0.48pt] (1,-4) circle (0.3);
    \draw[black,line width=0.48pt] (0.788,-3.788) -- (1.212,-4.212) (0.788,-4.212) -- (1.212,-3.788);
    \filldraw[fill=kgray,draw=black,line width=0.48pt] (2,-4) circle (0.3);
    \draw[black,line width=0.48pt] (1.788,-3.788) -- (2.212,-4.212) (1.788,-4.212) -- (2.212,-3.788);
    \filldraw[fill=white,draw=black,line width=0.48pt] (3,-4) circle (0.3);
    \filldraw[fill=white,draw=black,line width=0.48pt] (4,-4) circle (0.3);
    \filldraw[fill=white,draw=black,line width=0.48pt] (5,-4) circle (0.3);
    \filldraw[fill=kgray,draw=black,line width=0.48pt] (6,-4) circle (0.3);
    \draw[black,line width=0.48pt] (5.788,-3.788) -- (6.212,-4.212) (5.788,-4.212) -- (6.212,-3.788);
    \filldraw[fill=white,draw=black,line width=0.48pt] (7,-4) circle (0.3);
    \filldraw[fill=white,draw=black,line width=0.48pt] (8,-4) circle (0.3);
    \filldraw[fill=white,draw=black,line width=0.48pt] (9,-4) circle (0.3);
    \filldraw[fill=white,draw=black,line width=0.48pt] (10,-4) circle (0.3);
    \filldraw[fill=white,draw=black,line width=0.48pt] (11,-4) circle (0.3);
    \filldraw[fill=white,draw=black,line width=0.48pt] (12,-4) circle (0.3);
    \filldraw[fill=white,draw=black,line width=0.48pt] (13,-4) circle (0.3);
    \filldraw[fill=white,draw=black,line width=0.48pt] (1,-5) circle (0.3);
    \filldraw[fill=kgray,draw=black,line width=0.48pt] (2,-5) circle (0.3);
    \draw[black,line width=0.48pt] (1.788,-4.788) -- (2.212,-5.212) (1.788,-5.212) -- (2.212,-4.788);
    \filldraw[fill=kgray,draw=black,line width=0.48pt] (3,-5) circle (0.3);
    \draw[black,line width=0.48pt] (2.788,-4.788) -- (3.212,-5.212) (2.788,-5.212) -- (3.212,-4.788);
    \filldraw[fill=kgray,draw=black,line width=0.48pt] (4,-5) circle (0.3);
    \draw[black,line width=0.48pt] (3.788,-4.788) -- (4.212,-5.212) (3.788,-5.212) -- (4.212,-4.788);
    \filldraw[fill=kgray,draw=black,line width=0.48pt] (5,-5) circle (0.3);
    \draw[black,line width=0.48pt] (4.788,-4.788) -- (5.212,-5.212) (4.788,-5.212) -- (5.212,-4.788);
    \filldraw[fill=white,draw=black,line width=0.48pt] (6,-5) circle (0.3);
    \filldraw[fill=white,draw=black,line width=0.48pt] (7,-5) circle (0.3);
    \filldraw[fill=white,draw=black,line width=0.48pt] (8,-5) circle (0.3);
    \filldraw[fill=white,draw=black,line width=0.48pt] (9,-5) circle (0.3);
    \filldraw[fill=white,draw=black,line width=0.48pt] (10,-5) circle (0.3);
    \filldraw[fill=white,draw=black,line width=0.48pt] (11,-5) circle (0.3);
    \filldraw[fill=white,draw=black,line width=0.48pt] (12,-5) circle (0.3);
    \filldraw[fill=white,draw=black,line width=0.48pt] (13,-5) circle (0.3);
    \node[lab] at (1,-0.45) {1};
    \node[lab] at (2,-0.45) {2};
    \node[lab] at (3,-0.45) {3};
    \node[lab] at (4,-0.45) {4};
    \node[lab] at (5,-0.45) {5};
    \node[lab] at (6,-0.45) {6};
    \node[lab] at (7,-0.45) {7};
    \node[lab] at (8,-0.45) {8};
    \node[lab] at (9,-0.45) {9};
    \node[lab] at (10,-0.45) {10};
    \node[lab] at (11,-0.45) {11};
    \node[lab] at (12,-0.45) {12};
    \node[lab] at (13,-0.45) {13};
    \node[lab] at (0.45,-1) {1};
    \node[lab] at (0.45,-2) {2};
    \node[lab] at (0.45,-3) {3};
    \node[lab] at (0.45,-4) {4};
    \node[lab] at (0.45,-5) {5};
    \node[capt,anchor=north] at (7.00,-5.45) {(a)};
  \end{scope}
\end{tikzpicture}

\par\medskip
\begin{tikzpicture}[x=0.8cm,y=0.8cm,baseline=(current bounding box.north)]
  \begin{scope}[xshift=0.000cm]
    \draw[black,line width=0.32pt,dash pattern=on 1.20pt off 1.20pt] (1,-1) -- (16,-1) (1,-2) -- (16,-2) (1,-3) -- (16,-3) (1,-4) -- (16,-4) (1,-5) -- (16,-5) (1,-6) -- (16,-6) (1,-1) -- (1,-6) (2,-1) -- (2,-6) (3,-1) -- (3,-6) (4,-1) -- (4,-6) (5,-1) -- (5,-6) (6,-1) -- (6,-6) (7,-1) -- (7,-6) (8,-1) -- (8,-6) (9,-1) -- (9,-6) (10,-1) -- (10,-6) (11,-1) -- (11,-6) (12,-1) -- (12,-6) (13,-1) -- (13,-6) (14,-1) -- (14,-6) (15,-1) -- (15,-6) (16,-1) -- (16,-6);
    \filldraw[fill=white,draw=black,line width=0.48pt] (1,-1) circle (0.3);
    \filldraw[fill=kgray,draw=black,line width=0.48pt] (2,-1) circle (0.3);
    \draw[black,line width=0.48pt] (1.788,-0.788) -- (2.212,-1.212) (1.788,-1.212) -- (2.212,-0.788);
    \filldraw[fill=white,draw=black,line width=0.48pt] (3,-1) circle (0.3);
    \filldraw[fill=white,draw=black,line width=0.48pt] (4,-1) circle (0.3);
    \filldraw[fill=white,draw=black,line width=0.48pt] (5,-1) circle (0.3);
    \filldraw[fill=kgray,draw=black,line width=0.48pt] (6,-1) circle (0.3);
    \draw[black,line width=0.48pt] (5.788,-0.788) -- (6.212,-1.212) (5.788,-1.212) -- (6.212,-0.788);
    \filldraw[fill=white,draw=black,line width=0.48pt] (7,-1) circle (0.3);
    \filldraw[fill=white,draw=black,line width=0.48pt] (8,-1) circle (0.3);
    \filldraw[fill=white,draw=black,line width=0.48pt] (9,-1) circle (0.3);
    \filldraw[fill=white,draw=black,line width=0.48pt] (10,-1) circle (0.3);
    \filldraw[fill=white,draw=black,line width=0.48pt] (11,-1) circle (0.3);
    \filldraw[fill=white,draw=black,line width=0.48pt] (12,-1) circle (0.3);
    \filldraw[fill=white,draw=black,line width=0.48pt] (13,-1) circle (0.3);
    \filldraw[fill=white,draw=black,line width=0.48pt] (14,-1) circle (0.3);
    \filldraw[fill=white,draw=black,line width=0.48pt] (15,-1) circle (0.3);
    \filldraw[fill=kgray,draw=black,line width=0.48pt] (16,-1) circle (0.3);
    \draw[black,line width=0.48pt] (15.788,-0.788) -- (16.212,-1.212) (15.788,-1.212) -- (16.212,-0.788);
    \filldraw[fill=white,draw=black,line width=0.48pt] (1,-2) circle (0.3);
    \filldraw[fill=white,draw=black,line width=0.48pt] (2,-2) circle (0.3);
    \filldraw[fill=kgray,draw=black,line width=0.48pt] (3,-2) circle (0.3);
    \draw[black,line width=0.48pt] (2.788,-1.788) -- (3.212,-2.212) (2.788,-2.212) -- (3.212,-1.788);
    \filldraw[fill=white,draw=black,line width=0.48pt] (4,-2) circle (0.3);
    \filldraw[fill=white,draw=black,line width=0.48pt] (5,-2) circle (0.3);
    \filldraw[fill=white,draw=black,line width=0.48pt] (6,-2) circle (0.3);
    \filldraw[fill=white,draw=black,line width=0.48pt] (7,-2) circle (0.3);
    \filldraw[fill=white,draw=black,line width=0.48pt] (8,-2) circle (0.3);
    \filldraw[fill=white,draw=black,line width=0.48pt] (9,-2) circle (0.3);
    \filldraw[fill=white,draw=black,line width=0.48pt] (10,-2) circle (0.3);
    \filldraw[fill=white,draw=black,line width=0.48pt] (11,-2) circle (0.3);
    \filldraw[fill=white,draw=black,line width=0.48pt] (12,-2) circle (0.3);
    \filldraw[fill=white,draw=black,line width=0.48pt] (13,-2) circle (0.3);
    \filldraw[fill=white,draw=black,line width=0.48pt] (14,-2) circle (0.3);
    \filldraw[fill=kgray,draw=black,line width=0.48pt] (15,-2) circle (0.3);
    \draw[black,line width=0.48pt] (14.788,-1.788) -- (15.212,-2.212) (14.788,-2.212) -- (15.212,-1.788);
    \filldraw[fill=kgray,draw=black,line width=0.48pt] (16,-2) circle (0.3);
    \draw[black,line width=0.48pt] (15.788,-1.788) -- (16.212,-2.212) (15.788,-2.212) -- (16.212,-1.788);
    \filldraw[fill=kgray,draw=black,line width=0.48pt] (1,-3) circle (0.3);
    \draw[black,line width=0.48pt] (0.788,-2.788) -- (1.212,-3.212) (0.788,-3.212) -- (1.212,-2.788);
    \filldraw[fill=white,draw=black,line width=0.48pt] (2,-3) circle (0.3);
    \filldraw[fill=white,draw=black,line width=0.48pt] (3,-3) circle (0.3);
    \filldraw[fill=white,draw=black,line width=0.48pt] (4,-3) circle (0.3);
    \filldraw[fill=white,draw=black,line width=0.48pt] (5,-3) circle (0.3);
    \filldraw[fill=white,draw=black,line width=0.48pt] (6,-3) circle (0.3);
    \filldraw[fill=white,draw=black,line width=0.48pt] (7,-3) circle (0.3);
    \filldraw[fill=white,draw=black,line width=0.48pt] (8,-3) circle (0.3);
    \filldraw[fill=white,draw=black,line width=0.48pt] (9,-3) circle (0.3);
    \filldraw[fill=white,draw=black,line width=0.48pt] (10,-3) circle (0.3);
    \filldraw[fill=white,draw=black,line width=0.48pt] (11,-3) circle (0.3);
    \filldraw[fill=white,draw=black,line width=0.48pt] (12,-3) circle (0.3);
    \filldraw[fill=white,draw=black,line width=0.48pt] (13,-3) circle (0.3);
    \filldraw[fill=kgray,draw=black,line width=0.48pt] (14,-3) circle (0.3);
    \draw[black,line width=0.48pt] (13.788,-2.788) -- (14.212,-3.212) (13.788,-3.212) -- (14.212,-2.788);
    \filldraw[fill=kgray,draw=black,line width=0.48pt] (15,-3) circle (0.3);
    \draw[black,line width=0.48pt] (14.788,-2.788) -- (15.212,-3.212) (14.788,-3.212) -- (15.212,-2.788);
    \filldraw[fill=white,draw=black,line width=0.48pt] (16,-3) circle (0.3);
    \filldraw[fill=white,draw=black,line width=0.48pt] (1,-4) circle (0.3);
    \filldraw[fill=kgray,draw=black,line width=0.48pt] (2,-4) circle (0.3);
    \draw[black,line width=0.48pt] (1.788,-3.788) -- (2.212,-4.212) (1.788,-4.212) -- (2.212,-3.788);
    \filldraw[fill=white,draw=black,line width=0.48pt] (3,-4) circle (0.3);
    \filldraw[fill=white,draw=black,line width=0.48pt] (4,-4) circle (0.3);
    \filldraw[fill=white,draw=black,line width=0.48pt] (5,-4) circle (0.3);
    \filldraw[fill=white,draw=black,line width=0.48pt] (6,-4) circle (0.3);
    \filldraw[fill=white,draw=black,line width=0.48pt] (7,-4) circle (0.3);
    \filldraw[fill=white,draw=black,line width=0.48pt] (8,-4) circle (0.3);
    \filldraw[fill=white,draw=black,line width=0.48pt] (9,-4) circle (0.3);
    \filldraw[fill=kgray,draw=black,line width=0.48pt] (10,-4) circle (0.3);
    \draw[black,line width=0.48pt] (9.788,-3.788) -- (10.212,-4.212) (9.788,-4.212) -- (10.212,-3.788);
    \filldraw[fill=kgray,draw=black,line width=0.48pt] (11,-4) circle (0.3);
    \draw[black,line width=0.48pt] (10.788,-3.788) -- (11.212,-4.212) (10.788,-4.212) -- (11.212,-3.788);
    \filldraw[fill=kgray,draw=black,line width=0.48pt] (12,-4) circle (0.3);
    \draw[black,line width=0.48pt] (11.788,-3.788) -- (12.212,-4.212) (11.788,-4.212) -- (12.212,-3.788);
    \filldraw[fill=kgray,draw=black,line width=0.48pt] (13,-4) circle (0.3);
    \draw[black,line width=0.48pt] (12.788,-3.788) -- (13.212,-4.212) (12.788,-4.212) -- (13.212,-3.788);
    \filldraw[fill=kgray,draw=black,line width=0.48pt] (14,-4) circle (0.3);
    \draw[black,line width=0.48pt] (13.788,-3.788) -- (14.212,-4.212) (13.788,-4.212) -- (14.212,-3.788);
    \filldraw[fill=white,draw=black,line width=0.48pt] (15,-4) circle (0.3);
    \filldraw[fill=white,draw=black,line width=0.48pt] (16,-4) circle (0.3);
    \filldraw[fill=white,draw=black,line width=0.48pt] (1,-5) circle (0.3);
    \filldraw[fill=white,draw=black,line width=0.48pt] (2,-5) circle (0.3);
    \filldraw[fill=kgray,draw=black,line width=0.48pt] (3,-5) circle (0.3);
    \draw[black,line width=0.48pt] (2.788,-4.788) -- (3.212,-5.212) (2.788,-5.212) -- (3.212,-4.788);
    \filldraw[fill=white,draw=black,line width=0.48pt] (4,-5) circle (0.3);
    \filldraw[fill=white,draw=black,line width=0.48pt] (5,-5) circle (0.3);
    \filldraw[fill=white,draw=black,line width=0.48pt] (6,-5) circle (0.3);
    \filldraw[fill=white,draw=black,line width=0.48pt] (7,-5) circle (0.3);
    \filldraw[fill=white,draw=black,line width=0.48pt] (8,-5) circle (0.3);
    \filldraw[fill=kgray,draw=black,line width=0.48pt] (9,-5) circle (0.3);
    \draw[black,line width=0.48pt] (8.788,-4.788) -- (9.212,-5.212) (8.788,-5.212) -- (9.212,-4.788);
    \filldraw[fill=kgray,draw=black,line width=0.48pt] (10,-5) circle (0.3);
    \draw[black,line width=0.48pt] (9.788,-4.788) -- (10.212,-5.212) (9.788,-5.212) -- (10.212,-4.788);
    \filldraw[fill=white,draw=black,line width=0.48pt] (11,-5) circle (0.3);
    \filldraw[fill=white,draw=black,line width=0.48pt] (12,-5) circle (0.3);
    \filldraw[fill=white,draw=black,line width=0.48pt] (13,-5) circle (0.3);
    \filldraw[fill=white,draw=black,line width=0.48pt] (14,-5) circle (0.3);
    \filldraw[fill=white,draw=black,line width=0.48pt] (15,-5) circle (0.3);
    \filldraw[fill=white,draw=black,line width=0.48pt] (16,-5) circle (0.3);
    \filldraw[fill=kgray,draw=black,line width=0.48pt] (1,-6) circle (0.3);
    \draw[black,line width=0.48pt] (0.788,-5.788) -- (1.212,-6.212) (0.788,-6.212) -- (1.212,-5.788);
    \filldraw[fill=white,draw=black,line width=0.48pt] (2,-6) circle (0.3);
    \filldraw[fill=white,draw=black,line width=0.48pt] (3,-6) circle (0.3);
    \filldraw[fill=kgray,draw=black,line width=0.48pt] (4,-6) circle (0.3);
    \draw[black,line width=0.48pt] (3.788,-5.788) -- (4.212,-6.212) (3.788,-6.212) -- (4.212,-5.788);
    \filldraw[fill=kgray,draw=black,line width=0.48pt] (5,-6) circle (0.3);
    \draw[black,line width=0.48pt] (4.788,-5.788) -- (5.212,-6.212) (4.788,-6.212) -- (5.212,-5.788);
    \filldraw[fill=kgray,draw=black,line width=0.48pt] (6,-6) circle (0.3);
    \draw[black,line width=0.48pt] (5.788,-5.788) -- (6.212,-6.212) (5.788,-6.212) -- (6.212,-5.788);
    \filldraw[fill=kgray,draw=black,line width=0.48pt] (7,-6) circle (0.3);
    \draw[black,line width=0.48pt] (6.788,-5.788) -- (7.212,-6.212) (6.788,-6.212) -- (7.212,-5.788);
    \filldraw[fill=kgray,draw=black,line width=0.48pt] (8,-6) circle (0.3);
    \draw[black,line width=0.48pt] (7.788,-5.788) -- (8.212,-6.212) (7.788,-6.212) -- (8.212,-5.788);
    \filldraw[fill=kgray,draw=black,line width=0.48pt] (9,-6) circle (0.3);
    \draw[black,line width=0.48pt] (8.788,-5.788) -- (9.212,-6.212) (8.788,-6.212) -- (9.212,-5.788);
    \filldraw[fill=white,draw=black,line width=0.48pt] (10,-6) circle (0.3);
    \filldraw[fill=white,draw=black,line width=0.48pt] (11,-6) circle (0.3);
    \filldraw[fill=white,draw=black,line width=0.48pt] (12,-6) circle (0.3);
    \filldraw[fill=white,draw=black,line width=0.48pt] (13,-6) circle (0.3);
    \filldraw[fill=white,draw=black,line width=0.48pt] (14,-6) circle (0.3);
    \filldraw[fill=white,draw=black,line width=0.48pt] (15,-6) circle (0.3);
    \filldraw[fill=white,draw=black,line width=0.48pt] (16,-6) circle (0.3);
    \node[lab] at (1,-0.45) {1};
    \node[lab] at (2,-0.45) {2};
    \node[lab] at (3,-0.45) {3};
    \node[lab] at (4,-0.45) {4};
    \node[lab] at (5,-0.45) {5};
    \node[lab] at (6,-0.45) {6};
    \node[lab] at (7,-0.45) {7};
    \node[lab] at (8,-0.45) {8};
    \node[lab] at (9,-0.45) {9};
    \node[lab] at (10,-0.45) {10};
    \node[lab] at (11,-0.45) {11};
    \node[lab] at (12,-0.45) {12};
    \node[lab] at (13,-0.45) {13};
    \node[lab] at (14,-0.45) {14};
    \node[lab] at (15,-0.45) {15};
    \node[lab] at (16,-0.45) {16};
    \node[lab] at (0.45,-1) {1};
    \node[lab] at (0.45,-2) {2};
    \node[lab] at (0.45,-3) {3};
    \node[lab] at (0.45,-4) {4};
    \node[lab] at (0.45,-5) {5};
    \node[lab] at (0.45,-6) {6};
    \node[capt,anchor=north] at (8.50,-6.45) {(b)};
  \end{scope}
\end{tikzpicture}

\caption{(a) The $5\times13$ test board of $S_8$. (b) The $6\times16$ blackout $S_{16}$.}
\label{fig:boards}
\end{figure}

\begin{theorem}[five rows, (a) and (c) computer-assisted, (b) checked in Lean]\label{thm:fiverow}
\leavevmode
\begin{enumerate}[(a)]
\item We have $f(5,m)=m+5$ for $m=5$, $6$, and $7$.
\item We have $f(5,m)=m+6$ for every $m\ge8$, attained by $S_8$ followed by $m-8$ columns black in
row $1$ only.
\item The $5\times8$ grid has exactly four maximum blackouts, the images of $S_8$ under row
reflection and column reversal.
\end{enumerate}
\end{theorem}

\begin{proof}
(b) The upper bound is Theorem~\ref{thm:rowbound}.
For the construction, $W_5=5$, so the collar theorem with right padding
in row $1$ reduces the whole family from $S_8$ to Figure~\ref{fig:boards}(a).
Three enumerators and Lean verify that the $954$ rectangles of this
test board have different presentations.
(a) At $m=5$, use Corollary~\ref{cor:small-m} and the construction $S_0$
in Theorem~\ref{thm:fivefive}. At $m=6,7$, the constructions are
$15,\allowbreak 12,\allowbreak 23,\allowbreak 3,\allowbreak 34,\allowbreak 45$ and
$1,\allowbreak 12,\allowbreak 23,\allowbreak 3,\allowbreak 34,\allowbreak 45,\allowbreak 15$,
recorded in \nolinkurl{jis/data/small_cells_sat_5x6.json} and
\nolinkurl{jis/data/small_cells_sat_5x7_b12.json}. They give the lower
bounds. Checked SAT certificates rule out $m+6$ points
(Section~\ref{sec:verify}).
(c) Two SAT enumerations find the four images. A separate column search
checks all $10\,(8)_6\,5^2=5{,}040{,}000$ equality candidates from
Theorem~\ref{thm:equality} and finds the same four.
\end{proof}

Theorem~\ref{thm:A} gives $f(h,m)\le h+m$ whenever a side is at most
four, and on $5\times5$, $5\times6$, $5\times7$, and $6\times6$.
Up to transposition, these exhaust the grids of area below $40$.
Thus $5\times8$ is the first failure of the guess $f(h,m)=h+m$.

\begin{remark}[five-row counts, computer-assisted]\label{rem:fivecounts}
For $m=13,\dots,23$ the counts $a_5(m)$ are
\begin{gather*}
193{,}406,\quad 1{,}346{,}758,\quad 8{,}650{,}416,\quad 52{,}128{,}242,\\
299{,}349{,}210,\quad 1{,}653{,}637{,}492,\quad 8{,}828{,}301{,}226,\quad 45{,}743{,}165{,}894,\\
230{,}934{,}909{,}328,\quad 1{,}139{,}486{,}126{,}866,\quad 5{,}510{,}071{,}794{,}736.
\end{gather*}
Each count agrees in two independent transfer calculations using the
equality structure of Theorem~\ref{thm:equality}.
% TODO(author): A-number once approved (the counts a_5(m) for m >= 8).
\end{remark}

\subsection{Six rows and small grids}

Let $S_{16}$ be the blackout of the $6\times16$ grid in Figure~\ref{fig:boards}(b). Its doubles,
in the columns $1,2,3,6,9,10,14,15,16$, carry the edges of $K_{3,3}$ between $\{1,3,5\}$ and
$\{2,4,6\}$.

\begin{theorem}[six rows, (a) computer-assisted, (b) checked in Lean]\label{thm:sixrow}
\leavevmode
\begin{enumerate}[(a)]
\item For $8\le m\le14$ we have $f(6,m)\le m+8$.
\item For every $m\ge15$ we have $f(6,m)=m+9$.
\end{enumerate}
\end{theorem}

\begin{proof}
(a) Corollary~\ref{cor:small-m} handles $m=8$.
For $9\le m\le14$, an optimum with $m+9$ points would have double
columns carrying $K_{3,3}$, by Theorem~\ref{thm:equality}.
There are ten possible row partitions. The structured encoding has
checked LRAT certificates excluding every one (Lemma~\ref{lem:enc}).
(b) Again the row bound gives the upper bound. For $m=15$, take
$56,\allowbreak 45,\allowbreak 34,\allowbreak 23,\allowbreak 1,\allowbreak 1,\allowbreak 14,\allowbreak 1,\allowbreak 1,\allowbreak 12,\allowbreak 36,\allowbreak 1,\allowbreak 25,\allowbreak 6,\allowbreak 16$
from \nolinkurl{jis/data/six_row_sat_w15_A.json}, which is valid with
$24$ points. For $m\ge16$, apply the collar theorem to $S_{16}$ in
Figure~\ref{fig:boards}(b). Since $W_6=7$, it suffices to add
$(17,1),\dots,(23,1)$. Three enumerators and Lean verify distinct
presentations for all $4{,}627$ rectangles of that test board.
\end{proof}

\begin{theorem}[small grids, computer-assisted]\label{thm:refute}
We have $f(6,6)=12$, $f(6,7)=13$, $f(7,7)=14$, and $f(6,8)=15$. The blackout
$15,\allowbreak 1,\allowbreak 123,\allowbreak 34,\allowbreak 4,\allowbreak 4,\allowbreak 456,\allowbreak 26$ of the $6\times8$ grid attains the last value.
\end{theorem}

\begin{proof}
The other constructions are $14,\allowbreak 1,\allowbreak 12,\allowbreak 2,\allowbreak 23,\allowbreak 3456$
on the $6\times6$ grid,
$16,\allowbreak 1,\allowbreak 1,\allowbreak 12,\allowbreak 2,\allowbreak 23,\allowbreak 3456$
on $6\times7$, and
$17,\allowbreak 1,\allowbreak 12,\allowbreak 2,\allowbreak 23,\allowbreak 3456,\allowbreak 67$
on $7\times7$. Their records are
\nolinkurl{jis/data/small_cells_sat_6x6.json},
\nolinkurl{jis/data/small_cells_sat_6x7.json}, and
\nolinkurl{jis/data/small_cells_sat_7x7.json}.
All are valid. For each grid, a checked SAT certificate excludes
one additional point.
\end{proof}

\begin{theorem}[the $5\times5$ grid]\label{thm:fivefive}
\leavevmode
\begin{enumerate}[(a)]
\item We have $f(5,5)=10$.
\item (Computer-assisted.) The $a_5(5)=8$ maximum blackouts are the images of $S_0$
below under the square's symmetries.
\end{enumerate}
\end{theorem}

\begin{proof}
(a) Corollary~\ref{cor:small-m} gives the upper bound. To attain it,
take $S_0$ to be $45,\allowbreak 34,\allowbreak 123,\allowbreak 1,\allowbreak 15$.
Its row-column graph is an $8$-cycle with two pendant edges, hence it is
axis-valid. Lemma~\ref{lem:leftmost} lists $30$ tilted rectangles, namely
nine unit diamonds, sixteen others from Lemma~\ref{lem:tilted-narrow},
and five spanning the grid. The fifteen weak rectangles show the pairs
below. After an oblique pair, the point following the bar is a white
corner of the axis-aligned rectangle through it.
\begin{center}\footnotesize
\begin{tabular}{@{}llll@{}}
$(1,2),(2,1)\mid(1,1)$ & $(1,3),(2,2)\mid(1,2)$ & $(3,4),(4,3)\mid(4,4)$ & $(2,5),(3,4)\mid(3,5)$ \\
$(2,2),(4,2)$ & $(1,2),(4,3)\mid(1,3)$ & $(4,4),(5,2)\mid(4,2)$ & $(3,5),(4,4)\mid(3,4)$ \\
$(4,3),(5,2)\mid(4,2)$ & $(3,4),(5,2)\mid(5,4)$ & $(2,5),(5,2)\mid(2,2)$ & $(4,5),(5,3)\mid(4,3)$ \\
$(1,3),(4,2)\mid(1,2)$ & $(2,5),(4,3)\mid(4,5)$ & & $(4,5),(5,4)\mid(4,4)$ \\
\end{tabular}
\end{center}
These pairs are all different. The other tilted rectangles show at least
three points, so Lemma~\ref{lem:threecorners} distinguishes them.
The listed white corners rule out axis-aligned rivals for every oblique
pair. The remaining pair $(2,2),(4,2)$ has no rival because
$B_2\cap B_4=\emptyset$. This proves validity.
(b) Three programs check all $44{,}640$ sets of $10$ points whose graph
has no four- or six-cycle. Only the images of $S_0$ are valid.
\end{proof}
